\section{Fazit}
\label{sec:fazit}

Die vorliegende Projektarbeit hat gezeigt, wie eine analytische Datenanwendung zur Auswertung von Insider‑Transaktionen aufgebaut werden kann. Ausgehend von öffentlich zugänglichen SEC‑Meldungen wurden mithilfe der FMP‑API die relevanten Daten abgerufen, validiert, angereichert, bewertet und schließlich visualisiert.

Die Trennung zwischen Raw‑ und Clean‑Store ermöglicht eine flexible Speicherung der Rohdaten sowie eine performante Abfrage der bereinigten Daten. Die implementierte Pipeline ist modular aufgebaut und lässt sich leicht erweitern. Die Streamlit‑App bietet einen intuitiven Zugang zu den Daten und erlaubt interaktive Analysen. Durch die Wahl von Python und gängigen Open‑Source‑Bibliotheken bleibt die Lösung reproduzierbar und portabel.

Die wichtigste Einschränkung liegt im heuristischen Charakter des Scorings. Ohne empirische Validierung können die Ergebnisse lediglich als Illustrationen der technischen Möglichkeiten verstanden werden. Um aussagekräftige Handlungsempfehlungen ableiten zu können, wären umfangreichere Datenanalysen und statistische Modelle erforderlich.

Insgesamt erfüllt die entwickelte Anwendung die Anforderungen des Moduls und liefert einen umfassenden Einblick in die Verarbeitung, Speicherung und Analyse von Insider‑Transaktionsdaten. Sie legt damit den Grundstein für weitergehende Forschung und Anwendungen im Bereich der Kapitalmarktanalyse.